% BHU PYQ -- AIHC & Archaeology: Cultural Heritage of Varanasi (Vocational / Minor) (2025-26 NEP)
\documentclass[11pt,a4paper]{article}
\usepackage[a4paper,top=2.5cm,bottom=2.5cm,left=2.8cm,right=2.8cm,headheight=14pt]{geometry}
\usepackage{amsmath,amssymb}
\usepackage{fontspec}
\setmainfont{Kohinoor Devanagari}
\usepackage{microtype,fancyhdr,enumitem,hyperref,lastpage,parskip}

\hypersetup{colorlinks=true,urlcolor=black,linkcolor=black,pdftitle={AIHMV31: Cultural Heritage of Varanasi -- BHU NEP 2025-26}}

\pagestyle{fancy}\fancyhf{}
\renewcommand{\headrulewidth}{0.4pt}\renewcommand{\footrulewidth}{0.4pt}
\fancyhead[L]{\small \textit{Banaras Hindu University}}
\fancyhead[R]{\small \textit{AIHMV31: Cultural Heritage of Varanasi (2025-26)}}
\fancyfoot[C]{\small Page \thepage\ of \pageref{LastPage}}

\newlist{parts}{enumerate}{2}
\setlist[parts,1]{label=(\alph*),leftmargin=*,itemsep=4pt,topsep=3pt}
\setlist[parts,2]{label=(\roman*),leftmargin=*,itemsep=2pt,topsep=2pt}

\newcommand{\pts}[1]{\hfill \textit{[#1]}}

\begin{document}

\begin{center}
  {\large\bfseries BANARAS HINDU UNIVERSITY}\\[2pt]
  {\normalsize B. A. (Hons.) Semester III Examination, 2025-26 (NEP)}\\[6pt]
  \rule{0.8\linewidth}{0.5pt}\\[6pt]
  {\large\bfseries ANCIENT INDIAN HISTORY, CULTURE \& ARCHAEOLOGY}\\[4pt]
  {\large\bfseries Paper Code: AIHMV-31 : Cultural Heritage of Varanasi (Minor / Vocational)}\\[4pt]
  {\normalsize Time: 2:30 Hours \hfill Full Marks: 70}\\[2pt]
  {\small REF NO. 2025-26/D-III/070}
\end{center}
\rule{\linewidth}{0.4pt}
\smallskip

\noindent \textit{Instructions: Attempt all questions of Sections A, B and C which carry 10, 05 and 02 marks each respectively. Write your Roll No.\ at the top immediately on receipt of this question paper.}

\smallskip
\noindent \textit{खण्ड अ, ब तथा स के सभी प्रश्नों के उत्तर दीजिये जिनके अंक क्रमशः 10, 05 तथा 02 हैं।}

\bigskip

\section*{SECTION -- A / खण्ड -- अ}
\noindent \textit{Note: Long Answer Type Questions. Answer approximately in 500 words.}\pts{3 $\times$ 10 = 30}

\noindent \textit{दीर्घ उत्तरीय प्रश्न। लगभग 500 शब्दों में उत्तर दीजिये।}

\begin{enumerate}[label=\textbf{\arabic*.},leftmargin=*,itemsep=12pt]

  \item Write an essay on the antiquity and historical background of Varanasi.\pts{10}
  
  वाराणसी की प्राचीनता एवं ऐतिहासिक पृष्ठभूमि पर एक निबन्ध लिखिये।

  \begin{center}\textbf{OR / अथवा}\end{center}

  Discuss the significance of Sarnath as a major center of Buddhism.\pts{10}
  
  बौद्ध धर्म के प्रमुख केंद्र के रूप में सारनाथ के महत्त्व की विवेचना कीजिये।

  \item Describe the major Ghats of Varanasi and explain their cultural and religious importance.\pts{10}
  
  वाराणसी के प्रमुख घाटों का वर्णन कीजिये तथा उनके सांस्कृतिक एवं धार्मिक महत्त्व को समझाइये।

  \begin{center}\textbf{OR / अथवा}\end{center}

  Evaluate the role of Varanasi in the development of Indian Classical Music and Performing Arts.\pts{10}
  
  भारतीय शास्त्रीय संगीत एवं प्रदर्शन कलाओं के विकास में वाराणसी की भूमिका का मूल्यांकन कीजिये।

  \item Write a detailed note on the traditional handicrafts and silk weaving industry (Banarasi Saree) of Varanasi.\pts{10}
  
  वाराणसी के पारम्परिक हस्तशिल्प एवं रेशम बुनाई उद्योग (बनारसी साड़ी) पर एक विस्तृत टिप्पणी लिखिये।

\end{enumerate}

\bigskip

\section*{SECTION -- B / खण्ड -- ब}
\noindent \textit{Note: Short Answer Type Questions. Answer approximately in 250 words.}\pts{4 $\times$ 5 = 20}

\noindent \textit{लघु उत्तरीय प्रश्न। लगभग 250 शब्दों में उत्तर दीजिये।}

\begin{enumerate}[label=\textbf{\arabic*.},leftmargin=*,start=4,itemsep=10pt]

  \item Write a short note on Kashi Vishwanath Temple and its cultural heritage.\pts{5}
  
  काशी विश्वनाथ मन्दिर एवं इसकी सांस्कृतिक विरासत पर एक संक्षिप्त टिप्पणी लिखिये।

  \begin{center}\textbf{OR / अथवा}\end{center}

  Discuss the Panchakroshi Yatra of Varanasi.\pts{5}
  
  वाराणसी की पंचक्रोशी यात्रा की विवेचना कीजिये।

  \item Throw light on the literary contributions of Munshi Premchand and Bharatendu Harishchandra to Varanasi.\pts{5}
  
  वाराणसी के प्रति मुंशी प्रेमचंद एवं भारतेन्दु हरिश्चंद्र के साहित्यिक योगदानों पर प्रकाश डालिये।

  \begin{center}\textbf{OR / अथवा}\end{center}

  Write about the archaeological importance of Rajghat excavations in Varanasi.\pts{5}
  
  वाराणसी में राजघाट उत्खनन के पुरातात्विक महत्त्व के विषय में लिखिये।

  \item Discuss the concept of `Kashi Labha' and the philosophical significance of death in Varanasi.\pts{5}
  
  `काशी लाभ' की अवधारणा एवं वाराणसी में मृत्यु के दार्शनिक महत्त्व की चर्चा कीजिये।

  \begin{center}\textbf{OR / अथवा}\end{center}

  Explain the origin of the names `Kashi', `Varanasi', and `Banaras'.\pts{5}
  
  `काशी', `वाराणसी' एवं `बनारस' नामों की उत्पत्ति की व्याख्या कीजिये।

\end{enumerate}

\bigskip

\section*{SECTION -- C / खण्ड -- स}
\noindent \textit{Note: Very Short Answer Type Questions. Answer in the briefest possible way.}\pts{10 $\times$ 2 = 20}

\noindent \textit{अति संक्षिप्त उत्तर दीजिये।}

\begin{enumerate}[label=\textbf{\arabic*.},leftmargin=*,start=7,itemsep=8pt]
  \item 
  \begin{parts}
    \item What is the meaning of Folk religion?\pts{2}
    
    लोक धर्म (Folk religion) से आप क्या समझते हैं?

    \item Write the names of any two ancient Jain temples located in Varanasi.\pts{2}
    
    वाराणसी में स्थित जैन धर्म से सम्बद्ध किन्हीं दो प्राचीन मन्दिरों के नाम लिखिये। (Bhadaini / Bhelupur Parshvanath)

    \item Name four Brahmanical scriptures that provide information about Varanasi.\pts{2}
    
    वाराणसी के विषय में जानकारी प्रदान करने वाले किन्हीं चार ब्राह्मण ग्रन्थों के नाम लिखिये। (Kashi Khanda of Skanda Purana, Matsya Purana, Padma Purana, Mahabharata)

    \item What is the meaning of ``Tirtha''?\pts{2}
    
    ``तीर्थ'' का क्या अर्थ है?

    \item What is a Vihara and what was its purpose?\pts{2}
    
    विहार क्या है और इसका क्या उद्देश्य था? (Monastic residence for Buddhist monks)

    \item Name three famous books written by Munshi Premchand.\pts{2}
    
    मुंशी प्रेमचंद द्वारा लिखित किन्हीं तीन प्रसिद्ध पुस्तकों के नाम लिखिये। (Godan, Gaban, Sevasadan)

    \item Name two Buddhist stupas known from the archaeological excavated site of Sarnath.\pts{2}
    
    सारनाथ के उत्खनित पुरास्थल से ज्ञात किन्हीं दो बौद्ध स्तूपों के नाम लिखिये। (Dhamek Stupa, Chaukhandi Stupa)

    \item Who is the author of ``Bharat Durdasha''?\pts{2}
    
    ``भारत दुर्दशा'' के नाटककार/लेखक कौन हैं? (Bharatendu Harishchandra)

    \item Which two rivers border the city of Varanasi?\pts{2}
    
    वाराणसी शहर किन दो नदियों (वरुणा एवं अस्सी) की सीमाओं के मध्य स्थित है?

    \item Who founded the Banaras Hindu University (BHU) in 1916?\pts{2}
    
    1916 ई० में बनारस हिन्दू विश्वविद्यालय की स्थापना किसने की थी? (Pandit Madan Mohan Malaviya)
  \end{parts}
\end{enumerate}

\end{document}
